\section{Phase markers}
\label{sec:markers}

\subsection{Structure factor and the AFM label}
\label{sec:sq}

To identify the antiferromagnetic phase we use the structure factor
\begin{equation}
S(q)=\frac{1}{L}\sum_{i,j}e^{iq(i-j)}\langle Z_iZ_j\rangle,
\end{equation}
where $\langle Z_iZ_j\rangle$ is the ground-state correlator and $L=8$
throughout this section.
The structure factor is the Fourier transform of spin correlations; $S(\pi)$ quantifies staggered correlations in the finite chain. The $S(q)$ curves at one representative point in each
of the three phases (labels trivial / topo / AFM, coordinates in the
dataset) on the shared $q$ grid (including $q=\pi$) are shown in
Fig.~\ref{fig:D03}.
The $q=\pi$ component is used as the AFM diagnostic throughout the
hardware-facing protocol. It is evaluated from $Z_iZ_j$ correlations obtained by computational-basis measurements. Hardware $p=1$ measurements at the three points (Fig.~\ref{fig:D03}(b)) track the simulated curves in Fig.~\ref{fig:D03}(a), with the AFM $\pi$ peak clearly resolved despite noise suppression.

\begin{figure}[htbp]
\centering
\includegraphics[width=\columnwidth]{{../../figs/exp02/exp02_D03}}
\caption{Structure factor $S(q)$ at three phase-representative points (one
per phase) on the shared $q$ grid: (a) simulator, (b) hardware $p=1$
measurements at the same points ($1024$ shots with readout mitigation),
tracking the simulated curves; the values at $q=\pi$ compare staggered
correlations at the representative points.
\label{fig:D03}}
\end{figure}
